\subsection{Vector subspaces and quotient spaces of a topological vector space; products of topological vector spaces; topological direct sums of subspaces}

Everything that has been said for topological modules (GT, III, § 6.6) applies in particular to topological vector spaces. If M is a vector subspace of a topological vector space E, the topology induced on M by that of E is compatible with the vector space structure of M, and the closure \(\overline{M}\) of M in E is a vector subspace of E. The quotient topology of that of E by M is compatible with the vector space structure of E/M.

If E is a topological vector space, the closure N of \(\{0\}\) in E (intersection of neighbourhoods of 0) is a closed vector subspace of E; the quotient vector subspace E/N, which is necessarily Hausdorff whether E is or not, is called the Hausdorff vector space associated with E.

Let \((\mathrm{E}_{\iota})_{\iota \in I}\) be a family of topological vector spaces over the same topological division ring K, and let E be the product vector space of the \(E_{\iota}\) . The product topology of the topologies of the \(E_{\iota}\) is compatible with the vector space structure of E. In the product space E, the subspace F, the direct sum of the \(E_{\iota}\) is everywhere dense (GT, III, § 2.9, prop. 25).

For certain types of topological vector spaces on the field \(\mathbf{R}\) or the field \(\mathbf{C}\) we define (in II, p.~29) a topology on the direct sum of a family \((\mathrm{E}_{\mathrm{t}})\) of topological vector spaces that is, in general, distinct from the topology induced by the product topology of the \(\mathrm{E}_{\mathrm{t}}\) .

Everything that has been said on the finite direct sums of stable subgroups of topological groups with operators (GT, III, §6.2) applies to topological vector spaces, replacing « stable subgroup » throughout by « vector subspace ».

Remark. --- Given a closed vector subspace M of a Hausdorff topological vector space E, it is not necessarily the case that there exists an (algebraic) complementary vector subspace to M that is closed in E (even if E is a normed space; cf.~IV, p.~55, exerc. 16 (c)); a fortiori there does not necessarily exist a topological complement of M in E (cf.~I, p.~26, exerc. 8). However we shall see in § 2 that when K is a non-discrete valued division ring, then every closed subspace M of E, with finite codimension, does have a topological complement in E (I, p.~14, prop. 3).

